% Generated draft: requires mathematical and source review.
\documentclass[11pt]{article}
\usepackage[margin=25mm]{geometry}
\usepackage{fontspec}
\setmainfont{Latin Modern Roman}
\usepackage{amsmath,amssymb,mathtools,unicode-math}
\setmathfont{Latin Modern Math}
\usepackage{enumitem,needspace,xurl,hyperref}
\setlength{\parindent}{0pt}
\setlength{\parskip}{5pt}
\setlength{\emergencystretch}{3em}
\newcommand{\N}{\mathbb{N}}
\newcommand{\Z}{\mathbb{Z}}
\newcommand{\Q}{\mathbb{Q}}
\newcommand{\R}{\mathbb{R}}
\newcommand{\C}{\mathbb{C}}
\newcommand{\F}{\mathbb{F}}
\newcommand{\A}{\mathbb{A}}
\newcommand{\PP}{\mathbb{P}}
\newcommand{\E}{\mathbb{E}}
\newcommand{\eps}{\varepsilon}
\newcommand{\dd}{\,\mathrm{d}}
\begin{document}
\section*{UnsolvedMath ID: 30004827}
\textbf{Source problem code:} catalogue entry
\subsection{Definitions and mathematical statement}
Let \(R\) denote the hyperfinite factor of type \(\mathrm{II}_1\) referred to in the source. A factor is a von Neumann algebra with center \(\C 1\); type \(\mathrm{II}_1\) means that it has a faithful normal tracial state \(\tau\), normalized by \(\tau(1)=1\), and no nonzero minimal projections. Hyperfiniteness means approximation by finite-dimensional unital \(*\)-subalgebras whose union is weakly dense. The source does not explicitly specify a separability convention for \(R\).

A \(C^*\)-tensor category is a monoidal \(C^*\)-category with compatible \(*\)-preserving tensor product. It is rigid if every object has duality data satisfying the conjugate equations. The supplied excerpt does not specify size restrictions, countability, semisimplicity, or a simple tensor unit; these hypotheses must be recovered from the report's conventions before fixing the scope of the question.

For a rigid \(C^*\)-tensor category \(\mathcal C\), write \(\mathscr A_{\mathrm{free}}(\mathcal C,R)\) provisionally for the collection of free actions intended by the source, and write \(\sim_{\mathrm{cc}}\) for its intended notion of cocycle conjugacy. These symbols are placeholders: the supplied record does not define an action, freeness, or cocycle conjugacy, and the original excerpt puts quotation marks around ``act'' and ``cocycle conjugacy.'' Consequently, the following is a conditional formulation, pending those definitions.

\begin{enumerate}
\item Does every rigid \(C^*\)-tensor category \(\mathcal C\), within the intended scope, satisfy
\[
\mathscr A_{\mathrm{free}}(\mathcal C,R)\ne\varnothing?
\]
\item For every nonamenable rigid \(C^*\)-tensor category \(\mathcal C\) within that scope, does \(\mathscr A_{\mathrm{free}}(\mathcal C,R)\) have infinitely many equivalence classes under \(\sim_{\mathrm{cc}}\)? Here nonamenability must also be interpreted according to the report's conventions. The universal quantifier records the natural reading of ``in the nonamenable case''; the excerpt does not make that quantifier explicit.
\item What are the answers to these existence and multiplicity questions for the Temperley--Lieb--Jones categories at index \(I>4\)? The excerpt does not define the index convention or identify a more specific target for this last question.
\end{enumerate}

The supplied dated curation labels the problem partially solved and reports Rohlin and uniqueness results for amenable rigid \(C^*\)-tensor category actions on properly infinite factors, while describing existence and classification on the hyperfinite \(\mathrm{II}_1\) factor beyond amenability as unresolved. This is supplied background, not an independently verified theorem or current-status assessment; its properly infinite hypothesis must be retained.
\subsection{Short English statement}
The source asks whether every rigid \(C^*\)-tensor category admits a free action on the hyperfinite \(\mathrm{II}_1\) factor, whether each nonamenable category has infinitely many such actions up to cocycle conjugacy, and what happens for Temperley--Lieb--Jones categories of index greater than \(4\). Precise action and equivalence conventions, category hypotheses, and the scope of the final question remain unresolved in the supplied excerpt.
\subsection{Sources}
\begin{itemize}
\item Supplied local UnsolvedMath registry record, permanent ID \(30004827\), source problem code \texttt{catalogue entry}, titled ``Free Actions of Rigid Tensor Categories on the Hyperfinite Factor,'' in the Oberwolfach Reports Open Problems collection. No individual proposer is supplied; the registry gives proposed year \(2021\).
\item Supplied report citation: \emph{Quantum Groups -- Algebra, Analysis and Category Theory} (2022), Oberwolfach Reports, pp.~2397--2458, DOI \texttt{10.4171/owr/2021/44}. The original excerpt identifies p.~2405. \url{https://doi.org/10.4171/owr/2021/44}. This is also the supplied literature reference labeled ``Oberwolfach Report 44.''
\item Supplied literature reference labeled ``Tomatsu on centrally free actions'': \url{https://arxiv.org/abs/1812.04222}. The record supplies no further bibliographic details or exact theorem statement.
\item Literature assessment date: \texttt{2026-08-22}. The partially solved label and the assessment concerning amenable categories, properly infinite factors, and unresolved questions beyond amenability are supplied dated curation, not independently verified. The background explicitly requests independent verification of time-sensitive status and cited claims.
\item The registry's supplied statement status is \texttt{corrected\_verified}; its verification note says that the question was isolated from surrounding report prose and notation normalized. This supplied editorial label does not establish mathematical correctness or resolve missing definitions.
\end{itemize}
\end{document}
